\documentclass[11pt,a4paper]{article}
\usepackage[utf8]{inputenc}
\usepackage[T1]{fontenc}
\usepackage{tgheros}
\renewcommand{\familydefault}{\sfdefault}
\usepackage[varqu,scaled=0.92]{zi4}
\usepackage{microtype}
\DisableLigatures[-,<,>]{encoding=T1,family=tt*}
\usepackage{needspace}
\usepackage[a4paper,top=2.5cm,bottom=2.5cm,left=2.5cm,right=2.2cm]{geometry}
\usepackage{graphicx}
\graphicspath{{../../template/}{figures/}}
\usepackage{xcolor}
\usepackage{listings}
\usepackage{upquote}
\usepackage{titlesec}
\usepackage{tocloft}
\usepackage{enumitem}
\usepackage{booktabs}
\usepackage{array}
\usepackage{float}
\usepackage[hidelinks]{hyperref}

\setlength{\parindent}{0pt}
\setlength{\parskip}{0.55em}
\setlist{itemsep=0.3em,topsep=0.3em}
\pagestyle{plain}

\titleformat{\section}{\Large\bfseries}{\thesection}{0.9em}{}
\titleformat{\subsection}{\large\bfseries}{\thesubsection}{0.9em}{}
\titlespacing*{\section}{0pt}{1.6em}{0.7em}
\titlespacing*{\subsection}{0pt}{1.2em}{0.5em}

\renewcommand{\contentsname}{Cuprins}
\renewcommand{\cfttoctitlefont}{\LARGE\bfseries}
\renewcommand{\cftsecfont}{\bfseries}
\renewcommand{\cftsecpagefont}{\bfseries}
\renewcommand{\cftsubsecfont}{\small}
\renewcommand{\cftsubsecpagefont}{\small}
\setlength{\cftbeforesecskip}{0.6em}

\definecolor{codebg}{HTML}{F6F6F6}
\definecolor{coderule}{HTML}{8C8C8C}
\definecolor{codenum}{HTML}{8C8C8C}
\definecolor{codecomment}{HTML}{6A6A6A}

\lstdefinestyle{term}{
  basicstyle=\ttfamily\small,
  backgroundcolor=\color{codebg},
  frame=single,
  rulecolor=\color{coderule},
  framesep=4pt,
  numbers=left,
  numberstyle=\tiny\color{codenum},
  numbersep=9pt,
  xleftmargin=1.4em,
  breaklines=true,
  postbreak=\mbox{\textcolor{codenum}{$\hookrightarrow$}\space},
  columns=fullflexible,
  keepspaces=true,
  showstringspaces=false,
  upquote=true,
  aboveskip=0.9em,
  belowskip=0.9em
}
\lstdefinestyle{ccode}{
  style=term,
  language=C,
  keywordstyle=\bfseries,
  commentstyle=\color{codecomment},
  morekeywords={uint,uint64}
}
\lstset{style=term,captionpos=b}

\usepackage{caption}
\captionsetup[lstlisting]{labelsep=endash,justification=centering,singlelinecheck=false,font=small,skip=7pt}
\captionsetup[table]{labelsep=endash,font=small,skip=6pt}
\renewcommand{\lstlistingname}{Figure}
\AtBeginDocument{\counterwithin{lstlisting}{section}\counterwithin{table}{section}}

\newcommand{\plainsection}[1]{\section*{#1}\addcontentsline{toc}{section}{#1}}
\newcommand{\code}[1]{\texttt{#1}}

\newcommand{\utmtitlepage}{%
\begin{titlepage}
  \centering
  \includegraphics[width=4.3cm]{utm-logo.jpeg}\par
  \vspace{1.3cm}
  {\large\bfseries UNIVERSITATEA TEHNICĂ A MOLDOVEI\par}
  \vspace{0.25cm}
  Facultatea Calculatoare, Informatică și Microelectronică\par
  Departamentul Ingineria Software și Automatică\par
  \vspace{2.6cm}
  {\Large\bfseries LUCRARE DE LABORATOR NR.~\labnumber\par}
  \vspace{0.2cm}
  {\LARGE\bfseries \labtitle\par}
  \vspace{0.45cm}
  Disciplina: Criptografie și Securitate\par
  \ifx\variant\empty\else\vspace{0.2cm}Varianta \variant\par\fi
  \vspace{2.8cm}
  \begin{minipage}[t]{0.46\textwidth}
    \raggedright
    \textbf{A efectuat:}\\
    student gr. \studentgroup\\
    \studentname
  \end{minipage}\hfill
  \begin{minipage}[t]{0.46\textwidth}
    \raggedleft
    \textbf{A verificat:}\\
    conf. univ., dr. Aureliu \textsc{Zgureanu}
  \end{minipage}
  \vfill
  Chișinău — 2026\par
\end{titlepage}}

\usepackage{textcomp}
\usepackage{framed}
\newenvironment{cipherbox}{%
  \def\FrameCommand{\fboxsep=6pt\fboxrule=0.4pt\fcolorbox{coderule}{codebg}}%
  \MakeFramed{\advance\hsize-\width\FrameRestore}%
  \ttfamily\footnotesize\raggedright\frenchspacing\setlength{\parskip}{0pt}}%
  {\par\endMakeFramed}

\newcommand{\labnumber}{4}
\newcommand{\labtitle}{Block Ciphers: the DES Algorithm}
\newcommand{\studentname}{Daniel Nicorici}
\newcommand{\studentgroup}{FAF-243}
\newcommand{\variant}{16}
\lstset{inputencoding=utf8,extendedchars=true,literate=
  {Ă}{{\u{A}}}1 {Â}{{\^{A}}}1 {Î}{{\^{I}}}1 {Ș}{{\textcommabelow{S}}}1 {Ț}{{\textcommabelow{T}}}1
  {ă}{{\u{a}}}1 {â}{{\^{a}}}1 {î}{{\^{i}}}1 {ș}{{\textcommabelow{s}}}1 {ț}{{\textcommabelow{t}}}1}
\lstdefinestyle{py}{style=term,language=Python,keywordstyle=\bfseries,commentstyle=\color{codecomment}}
\begin{document}
\utmtitlepage
\tableofcontents
\newpage

\plainsection{Purpose of the work}

The purpose of this lab is to study the Data Encryption Standard (DES), the block cipher
standardised in FIPS 46, and to implement one element of it in a program that prints the tables it
uses and every intermediate step.

The task is chosen with the formula $nr\_sarcina = n \bmod 11$, where $n$ is the student's number
in the group register. My number is $n = 16$, and $16 \bmod 11 = 5$, so I solved task 2.5:
\textit{given K+, determine all 16 round keys $K_i$}. This is the DES key schedule. The program is
written in Python 3 with no external libraries. The user can type K+ or let the program generate
it at random.

\section{The DES algorithm}

Because the defence covers the whole algorithm, this section describes all of it before focusing
on the key schedule.

\subsection{Overview}

DES is a symmetric block cipher. It encrypts a 64-bit block of plaintext into a 64-bit block of
ciphertext with a 64-bit key, of which only 56 bits are used: every eighth bit (bits 8, 16, \dots,
64) is a parity bit. The same key decrypts. DES is a Feistel cipher with 16 rounds, so encryption
and decryption run the same steps, and only the order of the round keys differs.

\subsection{Encrypting one block}

\begin{enumerate}
  \item \textbf{Initial permutation.} The 64 bits of the block are reordered by the table IP and
  split into a left half $L_0$ and a right half $R_0$ of 32 bits each.
  \item \textbf{16 rounds.} For $i = 1, \dots, 16$:
  \[
    L_i = R_{i-1}, \qquad R_i = L_{i-1} \oplus f(R_{i-1}, K_i).
  \]
  Only the right half goes through the function $f$; the left half is just XORed with its result.
  \item \textbf{Final permutation.} After round 16 the halves are swapped, giving $R_{16}L_{16}$,
  and the inverse permutation $IP^{-1}$ produces the ciphertext.
\end{enumerate}

The round function $f(R_{i-1}, K_i)$ has four steps:
\begin{enumerate}
  \item \textbf{Expansion E.} The 32 bits of $R_{i-1}$ are expanded to 48 bits; 16 of them are
  repeated.
  \item \textbf{Key mixing.} The 48 bits are XORed with the round key: $K_i \oplus E(R_{i-1})$.
  \item \textbf{S-boxes.} The result is cut into eight 6-bit groups $B_1 \dots B_8$. For each group,
  the first and last bits form the row (0--3) and the four middle bits the column (0--15) of the
  table $S_j$, which gives a 4-bit number. The eight outputs together are 32 bits. The S-boxes are
  the only non-linear step of DES, and its security rests on them.
  \item \textbf{Permutation P.} The 32 bits are reordered by the table P.
\end{enumerate}

To decrypt, the same algorithm runs with the round keys in reverse order, $K_{16}, \dots, K_1$.

\subsection{The key schedule (task 2.5)}

The round keys come from the 64-bit key in three steps:
\begin{enumerate}
  \item \textbf{PC-1.} The permutation PC-1 keeps 56 of the 64 key bits (it drops the parity bits)
  and reorders them. The result is $K^+$, the input of my program.
  \item \textbf{Split and rotate.} $K^+$ is split into two 28-bit halves $C_0$ and $D_0$. For
  $i = 1, \dots, 16$, $C_i$ and $D_i$ are obtained from $C_{i-1}$ and $D_{i-1}$ by a cyclic shift to
  the left by 1 or 2 positions, according to Table~\ref{tab:shifts}.
  \item \textbf{PC-2.} The 56-bit concatenation $C_iD_i$ goes through PC-2
  (Table~\ref{tab:pc2}), which selects 48 of its bits in a fixed order. The result is the round key
  $K_i$: bit $j$ of $K_i$ is bit number PC-2$[j]$ of $C_iD_i$.
\end{enumerate}

\begin{table}[H]
\centering
\caption{Number of left shifts in each round}\label{tab:shifts}
\setlength{\tabcolsep}{4pt}
\small
\begin{tabular}{l*{16}{c}}
\toprule
Round $i$ & 1 & 2 & 3 & 4 & 5 & 6 & 7 & 8 & 9 & 10 & 11 & 12 & 13 & 14 & 15 & 16 \\
Shifts & 1 & 1 & 2 & 2 & 2 & 2 & 2 & 2 & 1 & 2 & 2 & 2 & 2 & 2 & 2 & 1 \\
\bottomrule
\end{tabular}
\end{table}

\begin{table}[H]
\centering
\caption{Permuted choice PC-2}\label{tab:pc2}
\small
\begin{tabular}{*{6}{c}}
\toprule
14 & 17 & 11 & 24 & 1 & 5 \\
3 & 28 & 15 & 6 & 21 & 10 \\
23 & 19 & 12 & 4 & 26 & 8 \\
16 & 7 & 27 & 20 & 13 & 2 \\
41 & 52 & 31 & 37 & 47 & 55 \\
30 & 40 & 51 & 45 & 33 & 48 \\
44 & 49 & 39 & 56 & 34 & 53 \\
46 & 42 & 50 & 36 & 29 & 32 \\
\bottomrule
\end{tabular}
\end{table}

\textbf{Observations:}
\begin{enumerate}
  \item \textbf{The shifts add up to 28.} Four rounds shift by 1 and twelve by 2, so
  $4 + 24 = 28$, a full turn of a 28-bit half. Therefore $C_{16} = C_0$ and $D_{16} = D_0$, and a
  decryptor can produce the keys in reverse order by rotating to the right.
  \item \textbf{PC-2 drops 8 bits.} Positions 9, 18, 22, 25, 35, 38, 43 and 54 never appear in
  PC-2, so 48 of the 56 bits reach each round key. The first 24 entries are all $\le 28$ and the
  last 24 all $> 28$: the first half of $K_i$ comes only from $C_i$ and the second half only from
  $D_i$.
  \item \textbf{Every key bit is used often.} Because the halves rotate, a different selection of
  bits reaches each round. Each of the 56 bits of $K^+$ appears in 12 to 15 of the 16 round keys,
  so every key bit affects many rounds.
\end{enumerate}

\section{Implementation}

The key schedule itself is the short function \code{round\_keys}. It splits $K^+$ into
\code{c} and \code{d}, and for every entry of \code{SHIFTS} it rotates both halves with string
slicing (\code{c[shift:] + c[:shift]}) and builds $K_i$ by reading $C_iD_i$ at the 48 positions
of \code{PC2}. The table positions start at 1, so the code subtracts 1 to index the string.

\begin{itemize}
  \item \code{parse\_kplus} accepts $K^+$ as 56 binary digits (spaces allowed) or as 14 hexadecimal
  digits, and refuses anything else with a message that states the allowed format.
  \item \code{show\_tables} prints the two tables the algorithm uses.
  \item \code{show\_schedule} prints $K^+$, $C_0$, $D_0$ and, for every round, the shift, $C_i$,
  $D_i$, $C_iD_i$ and $K_i$ in binary (grouped by 6 bits, the size of an S-box input) and in hex.
  For round 1 it also shows PC-2 at work, bit by bit.
  \item \code{main} is the menu: 1 reads $K^+$ from the user, 2 generates a random $K^+$ with the
  \code{secrets} module, 0 exits.
\end{itemize}

\lstinputlisting[style=py,caption={des\_keys.py},label={fig:code}]{../des_keys.py}

\section{Running the program}

The first run uses $K^+$ from the worked example of J.~Orlin Grabbe's tutorial \textit{The DES
Algorithm Illustrated}, so every value can be compared with a published source. It then tries an
invalid $K^+$ and an invalid menu choice.

\lstinputlisting[caption={Key schedule for the tutorial's K+},label={fig:run}]{run_example.txt}

The second run uses a random $K^+$. Rounds 2 to 15 are left out to save space.

\lstinputlisting[caption={Key schedule for a random K+ (excerpt)},label={fig:random}]{run_random.txt}

\textbf{Observations:}
\begin{enumerate}
  \item \textbf{The results match the tutorial.} $K_1 = $
  \code{000110 110000 001011 101111 111111 000111 000001 110010} and $K_{16} = $
  \code{110010 110011 110110 001011 000011 100001 011111 110101}, as in Grabbe's tutorial, and so
  do all the keys in between (checked automatically in section~\ref{sec:tests}).
  \item \textbf{$C_{16}D_{16} = K^+$.} In Figure~\ref{fig:run}, $C_{16}$ and $D_{16}$ are the same
  as $C_0$ and $D_0$, as predicted by the sum of the shifts.
  \item \textbf{Round 1 shows PC-2 directly.} The first row of PC-2 is 14, 17, 11, 24, 1, 5. The
  bits of $C_1D_1$ at those positions are 0, 0, 0, 1, 1, 0, which is exactly the first group
  \code{000110} of $K_1$.
\end{enumerate}

\section{Testing}\label{sec:tests}

The test file compares all 16 round keys, and $C_1$ and $D_1$, with the values printed in Grabbe's
tutorial. It also checks that $C_{16}D_{16} = K^+$, that the hexadecimal form of $K^+$ gives the
same result, which 8 positions PC-2 drops, that the four DES weak keys give 16 identical round keys,
and that invalid inputs are refused.

\lstinputlisting[style=py,caption={test\_des\_keys.py},label={fig:tests}]{../test_des_keys.py}

\begin{lstlisting}[caption={Running the tests},label={fig:testrun}]
$ python3 test_des_keys.py
all tests passed
\end{lstlisting}

The weak-key test shows a real flaw of the schedule. When $C_0$ and $D_0$ are each all zeros or all
ones, rotating them changes nothing, so every round uses the same key. Encrypting twice with such
a key gives back the plaintext, which is why these four keys must never be used.

\plainsection{Conclusions}

I implemented the DES key schedule: from the 56-bit $K^+$ the program produces the halves
$C_i$ and $D_i$ by cyclic left shifts and selects the 48-bit round keys $K_1, \dots, K_{16}$ with
PC-2, printing both tables and every intermediate value. All 16 keys match the worked example of a
published DES tutorial.

The key schedule is simple: rotations and fixed bit selections, with no arithmetic. Its job is to
give each round a different 48-bit view of the same 56-bit key, so that each key bit enters 12 to 15
of the rounds. Its structure also explains two properties of DES: the shifts sum to 28, so the
schedule returns to $K^+$ and can be run backwards for decryption, and keys with constant halves
are weak, because rotation cannot change them.

The real weakness of DES is not the schedule but the key length. 56 bits give $2^{56} \approx 7.2
\cdot 10^{16}$ keys, few enough that special hardware found a DES key by brute force in 1998. This
is why FIPS 46-3 added Triple DES and why DES has since been replaced by AES.
\end{document}
